% Ultima sezione del capitolo 3, dopo 3.2 (Gestione dell'incertezza sul moto delle persone).
% Non richiede rinumerazioni: 3.2.3 resta 3.2.3 e i rimandi esistenti continuano a valere.
% Riferimenti al codice: main.py, macchia_diffusione() e il blocco di aggregazione nel ciclo
% principale; costanti RAGGIO_BLOB_CELLE, DECADIMENTO_BLOB, COSTO_MASSIMO_PROBABILITA.

\subsection{Ottimizzazione della definizione delle regioni di rischio}
\label{sec:risoluzione-rischio}

La costruzione presentata in \S\ref{sec:aree-probabilita} lascia aperta una scelta che incide
sul risultato più di quanto la formulazione suggerisca, e che costituisce un compromesso fra
due grandezze in conflitto: la fedeltà geometrica con cui la regione di rischio è
rappresentata e il costo computazionale della ricerca del percorso. Entrambe dipendono dalla
risoluzione della griglia, ma in versi opposti, sicché adottare un'unica griglia per l'una e
per l'altra obbligherebbe a sacrificare la prima in favore della seconda o viceversa.

La ricerca del percorso preferisce celle grandi. Il costo della pianificazione cresce
approssimativamente con il cubo dell'inverso del passo, poiché dimezzare il lato della cella
quadruplica i nodi da espandere e raddoppia i campioni richiesti da ciascun controllo di
visibilità, che è l'operazione dominante del pianificatore. La regione di rischio richiede
invece celle piccole, essendo un oggetto geometrico: si allunga lungo la direzione di marcia,
si piega attorno agli spigoli e si restringe in corrispondenza dei varchi. Ciascuna di queste
proprietà è rappresentabile soltanto se la regione è descritta da un numero sufficiente di
celle; su una griglia dal passo confrontabile con l'ingombro di una persona l'anisotropia
degenera e la deformazione attorno a un angolo si riduce a un gradino.

L'ottimizzazione consiste nel disaccoppiare le due risoluzioni anziché mediarle. La macchia
viene propagata su una \textbf{sotto-griglia} di lato $\ell/k$, con $\ell$ il lato della cella
di pianificazione e $k$ intero, e il risultato è aggregato sulla griglia di pianificazione
prima che il pianificatore lo consumi. Il valore di $k$ è vincolato a essere multiplo del
rapporto fra le due griglie, cosicché l'aggregazione sia una divisione intera priva di resto e
nessuna sotto-cella ricada a cavallo di due celle di decisione.

Ne risulta un'architettura a due risoluzioni con ruoli distinti: una griglia uniforme e
relativamente grossolana per la \emph{decisione}, che ne determina il costo computazionale, e
una risoluzione locale e più fine per la \emph{rappresentazione del rischio}, che ne determina
la fedeltà geometrica. Perché tale separazione sia legittima devono però valere tre proprietà,
verificate nel seguito: che il raffinamento non alteri la regione che descrive, che
l'aggregazione non introduca a sua volta una dipendenza dalla risoluzione, e che il costo
aggiuntivo non ricada sulla componente che già domina il bilancio di calcolo.

\subsubsection{Invarianza della forma rispetto al raffinamento}

Il raffinamento deve modificare la fedeltà con cui la regione è descritta, non la regione
stessa: diversamente $k$ diventerebbe una seconda manopola sull'estensione del rischio, e non
sarebbe più possibile variarlo senza alterare al tempo stesso il comportamento del sistema.

Si è pertanto imposto che il raggio di propagazione, espresso in sotto-celle, scali con $k$, e
che il decadimento per sotto-cella valga $\delta^{1/k}$. La seconda condizione discende
direttamente dalla proprietà moltiplicativa stabilita in \S\ref{sec:aree-probabilita}: poiché
il peso di una cella è il prodotto dei decadimenti dei passi che vi conducono, elevare il
fattore all'inverso di $k$ compensa esattamente il numero $k$ volte maggiore di passi
necessari a percorrere la medesima distanza.

L'estensione fisica della regione e il decadimento per unità di lunghezza restano dunque
immutati al variare di $k$: raffinando si ottiene la stessa macchia descritta meglio, non una
macchia diversa. Il parametro governa esclusivamente la risoluzione della rappresentazione, e
può essere variato senza che alcun altro comportamento del sistema ne risulti alterato.

\subsubsection{Aggregazione per massimo}

A ciascuna cella di pianificazione è attribuito il \emph{massimo} fra i pesi delle sotto-celle
che la compongono, e non la loro somma né la loro media. Le due alternative sono state
scartate per ragioni distinte.

La somma renderebbe la penalità dipendente da $k$, poiché il numero di addendi cresce con
$k^2$: raffinare la rappresentazione equivarrebbe ad aumentare la prudenza del robot, due
effetti che devono restare distinti se si vuole poter variare l'uno senza l'altro. La media
diluirebbe invece un picco genuino nel volume della cella che lo contiene, attenuando proprio
la porzione di regione che il robot deve evitare: una sotto-cella a rischio elevato circondata
da sotto-celle libere produrrebbe una cella di costo modesto, e il pianificatore
l'attraverserebbe.

Il massimo conserva la lettura conservativa — il punto peggiore che il robot incontrerebbe
attraversando quella cella — ed è indipendente dalla risoluzione. Per la medesima ragione il
costo massimo associato al peso unitario è riscalato con il lato della cella: la penalità
accumulata percorrendo una data distanza non deve dipendere da come quella distanza è stata
suddivisa.

\subsubsection{Sostenibilità del costo computazionale}

Il beneficio di questa separazione non sarebbe interessante se il suo costo ricadesse sulla
ricerca del percorso, che domina il bilancio di calcolo. Non è così, per due ragioni.

La prima è che il pianificatore continua a operare sulla propria griglia: riceve una mappa di
costo indicizzata come prima dell'introduzione della sotto-griglia, e il numero di nodi che
espande non cambia. L'onere della ricerca — e con esso il controllo di visibilità che ne
costituisce la quasi totalità — resta quindi quello determinato dalla sola griglia di
pianificazione.

La seconda è che la propagazione raffinata è \emph{locale} per costruzione: si sviluppa
attorno a ciascuna persona rilevata, entro un raggio limitato e con interruzione al di sotto
di una soglia di peso trascurabile. Il suo costo è pertanto proporzionale al numero di persone
percepite, e non all'estensione dell'ambiente. Il raffinamento si paga dunque soltanto dove
serve, nell'intorno dei corpi osservati, mentre lo spazio libero — che costituisce la maggior
parte della mappa e sul quale non vi è alcun rischio da rappresentare — resta descritto alla
risoluzione grossolana che al pianificatore è sufficiente.

\subsubsection{Limite del beneficio e configurazione impiegata}

Il collasso della sotto-griglia sulla griglia di pianificazione merita di essere sottolineato,
poiché delimita l'ambito del beneficio ottenibile per questa via. Raffinare la
rappresentazione rende la regione di rischio geometricamente più fedele, ma non abbassa la
risoluzione alla quale il pianificatore decide, che resta quella della propria griglia.

Ne segue che l'ottimizzazione qui descritta migliora la qualità dell'informazione trasmessa e
non la finezza con cui essa può essere utilizzata. La distinzione risulterà determinante
nell'analisi condotta in \S\ref{sec:dove-si-arresta}, dove si mostra che il guadagno di
un'informazione più accurata si trasferisce alla decisione soltanto nella misura in cui supera
la risoluzione alla quale la decisione è presa: soglia che il raffinamento della sola
rappresentazione non sposta.

Va infine precisato, per la corretta lettura del capitolo successivo, che le campagne di
misura ivi riportate sono state eseguite con $k=1$, ossia con sotto-celle coincidenti con le
celle di pianificazione. Il meccanismo descritto appartiene dunque al sistema realizzato ma
non è esercitato da quei dati, e nessuna delle differenze misurate gli è attribuibile.
